\DocumentMetadata{
  pdfversion=2.0,
  pdfstandard=ua-2,
  lang=en-US,
  tagging=on,
  tagging-setup={math/setup=mathml-SE,tabsorder=structure}
}
\documentclass[11pt,letterpaper]{article}
\usepackage[margin=0.68in,includefoot]{geometry}
\usepackage{newcomputermodern}
\usepackage{amsmath,amscd,mathtools}
\usepackage{tikz,adjustbox}
\providecommand{\square}{\mdlgwhtsquare}
\usepackage[hidelinks]{hyperref}
\hypersetup{
  pdftitle={Numerical Analysis PhD Exam, January 2000},
  pdfauthor={University of Florida Department of Mathematics},
  pdfsubject={Graduate mathematics examination},
  pdfdisplaydoctitle=true,
  bookmarksopen=true
}
\setcounter{secnumdepth}{0}
\setlength{\parindent}{0pt}
\setlength{\parskip}{0.28em}
\setlength{\emergencystretch}{3em}
\setlength{\leftmargini}{2.15em}
\setlength{\leftmarginii}{2.65em}
\setlength{\leftmarginiii}{3em}
\renewcommand{\labelenumi}{\textbf{\theenumi.}}
\pagestyle{empty}
\raggedbottom
\allowdisplaybreaks
\definecolor{ProblemConcern}{HTML}{7A1F1F}
\newenvironment{problemconcern}{%
  \par\tagstructbegin{tag=Aside}\begingroup\color{ProblemConcern}
}{%
  \par\endgroup\tagstructend\nopagebreak[4]\vspace{0.12em}
}
\newenvironment{examproblems}[1]{%
  \begin{enumerate}%
  \setcounter{enumi}{#1}%
  \setlength{\topsep}{0.35em}%
  \setlength{\partopsep}{0pt}%
  \setlength{\itemsep}{0.48em}%
  \setlength{\parsep}{0.18em}%
}{\end{enumerate}}
\newenvironment{examsubparts}{%
  \begin{enumerate}%
  \setlength{\topsep}{0.18em}%
  \setlength{\partopsep}{0pt}%
  \setlength{\itemsep}{0.16em}%
  \setlength{\parsep}{0.1em}%
}{\end{enumerate}}
\newenvironment{examinstructions}{%
  \par\begingroup\itshape
}{%
  \par\endgroup\vspace{0.18em}
}
\makeatletter
\renewcommand\section{\@startsection{section}{1}{0pt}%
  {-1.25ex plus -.25ex minus -.1ex}%
  {0.55ex plus .1ex}%
  {\normalfont\large\bfseries}}
\renewcommand{\@maketitle}{%
  \newpage\null\vskip -1em%
  {\footnotesize\bfseries
    \href{https://www.ufl.edu/}{University of Florida}, \href{https://math.ufl.edu/}{Department of Mathematics}\par}%
  \vskip -0.5em%
  \tagstructbegin{tag=Title}%
    {\LARGE\bfseries\href{https://gma.math.ufl.edu/exams/numerical-analysis-phd/}{Numerical Analysis PhD Exam}, January 2000\par}%
  \tagstructend%
  \vskip 0.25em%
  \tagmcbegin{artifact}%
    \hrule height 1pt%
  \tagmcend%
  \vskip 1em%
}
\makeatother
\title{Numerical Analysis PhD Exam, January 2000}
\author{}
\date{}
\begin{document}
\maketitle
\thispagestyle{empty}
\begin{examinstructions}
Do any 8 of the following 10 problems.
\end{examinstructions}
\begin{examproblems}{0}
\item
This problem has four parts.
\begin{examsubparts}
\item[\textnormal{(1)}]
Show: If a matrix \(A \in \mathbb{R}^{n \times n}\) has eigenvalues \(\lambda_1,\ldots,\lambda_n\) such that \(\lambda_i \ne \lambda_j\) for all \(i \ne j\), then the corresponding eigenvectors are linearly independent.
\item[\textnormal{(2)}]
Show: If a symmetric matrix \(A \in \mathbb{R}^{n \times n}\) has eigenvalues \(\lambda_1,\ldots,\lambda_n\) such that \(\lambda_i \ne \lambda_j\) if \(i \ne j\), then the eigenvectors of~\(A\) are orthogonal.
\item[\textnormal{(3)}]
Show: If the matrix \(A \in \mathbb{R}^{n \times n}\) has~\(n\) independent eigenvectors, then the matrix can be diagonalized.
\item[\textnormal{(4)}]
Show: If \(A \in \mathbb{R}^{n \times n}\) is symmetric, then the matrix can be diagonalized by an orthogonal matrix.
\end{examsubparts}
\item
This problem has three parts.
\begin{examsubparts}
\item[\textnormal{(1)}]
Give a careful statement of the singular value decomposition for an arbitrary matrix \(A \in \mathbb{R}^{m \times n}\).
\item[\textnormal{(2)}]
Give a careful proof of the singular value decomposition.
\item[\textnormal{(3)}]
Show how the singular value decomposition for~\(A\) naturally yields a solution to the least squares problem \(\min_x \|Ax-y\|_2\).
\end{examsubparts}
\item
\begin{problemconcern}
\textbf{Warning: Suspected error.} The source calls~\(H\) an “orthogonal projection,” but \(H=I-2ww^{\mathsf{T}}\) is generally a Householder reflection (and an orthogonal matrix), not a projection. The source wording was retained because the intended replacement is not uniquely determined.
\end{problemconcern}
This problem has three parts.
\begin{examsubparts}
\item[\textnormal{(1)}]
Show: If~\(w\) is a unit vector in \(\mathbb{R}^n\), then the matrix \(H=I-2ww^t\) is an orthogonal projection.
\item[\textnormal{(2)}]
Show: If \(w=\frac{x-y}{\|x-y\|}\), \(\|x\|=\|y\|\), and \(H=I-2ww^t\), then \(Hx=y\).
\item[\textnormal{(3)}]
Given an arbitrary matrix \(A \in \mathbb{R}^{m \times n}\), construct an orthogonal matrix \(Q \in \mathbb{R}^{m \times m}\) and an upper triangular matrix \(R \in \mathbb{R}^{m \times n}\) such that \(A=QR\).
\end{examsubparts}
\item
This problem has three parts.
\begin{examsubparts}
\item[\textnormal{(1)}]
Show: If \(\lambda_n\) is the unique eigenvalue of \(A \in \mathbb{R}^{n \times n}\) with largest magnitude and \(x_0\) is chosen arbitrarily in \(\mathbb{R}^n\), then the iteration \(x_k=A^k x_0\) will almost surely converge to a multiple of \(v_n\), the eigenvector associated with \(\lambda_n\).
\item[\textnormal{(2)}]
If \(\lambda_k\) is a unique eigenvalue which does not have the largest magnitude from among the eigenvalues of~\(A\), then describe a method that will compute the eigenvector corresponding to \(\lambda_k\).
\item[\textnormal{(3)}]
Describe the \(QR\) iteration algorithm for computing the eigenvalues for a matrix \(A \in \mathbb{R}^{n \times n}\). Draw analogues between the \(QR\) method and the power method.
\end{examsubparts}
\item
BACKGROUND: It can be shown that for~\(\epsilon\) small there are differentiable functions \(x(\epsilon)\), and \(\lambda(\epsilon)\) such that \((A+\epsilon F)x(\epsilon)=\lambda(\epsilon)x(\epsilon)\), where \(\|F\|_2=1\). Let \(x=x(0)\) be a right eigenvector of~\(A\), and~\(y\) be a left eigenvector of~\(A\).
\begin{examsubparts}
\item[\textnormal{(1)}]
If \(A=D+F\), where \(A,D,F \in \mathbb{R}^{n \times n}\) and \(D=\operatorname{diag}(A)\), then state and prove a theorem relating the location of the eigenvalues of~\(A\) to the matrices~\(D\) and~\(F\).
\item[\textnormal{(2)}]
Show: If \(A,D,F \in \mathbb{R}^{n \times n}\), \(Ax'(0)+Fx=\lambda'(0)x+\lambda x'(0)\).
\item[\textnormal{(3)}]
Using part (2) above show that \(\left|\lambda'(0)\right|=\left\|\frac{y^tFx}{y^tx}\right\| \le \frac{1}{y^tx}\).
\item[\textnormal{(4)}]
Explain the implications of (3) to the stability of eigenvalues of symmetric matrices, and highly unsymmetric matrices.
\end{examsubparts}
\item
This problem has five parts.
\begin{examsubparts}
\item[\textnormal{(1)}]
Give a careful statement of the error formula for Lagrange interpolation.
\item[\textnormal{(2)}]
Give a careful proof of the error for formula for Lagrange interpolation.
\item[\textnormal{(3)}]
Give a careful definition of the Chebyshev Polynomials \(T_n(x)\).
\item[\textnormal{(4)}]
Give a careful statement of the theorem that shows that the roots of \(T_n(x)\) provide the optimal choice on \([-1,1]\) for Lagrange interpolation.
\item[\textnormal{(5)}]
For the function \(f(x)=\sin(\pi x)\) defined on \([-1,1]\) and polynomial interpolating function \(p_n(x)\) with interpolating points taken to be the roots of \(T_n(x)\), and tolerance \(\epsilon=0.00001\), find an integer~\(n\) such that the interpolating polynomial \(p_n(x)\) has the property that \(|p_n(x)-f(x)|<\epsilon\) for all \(x \in [-1,1]\).
\end{examsubparts}
\item
This problem has two parts.
\begin{examsubparts}
\item[\textnormal{(1)}]
Devise a 2nd order method for computing the cube root of a number.
\item[\textnormal{(2)}]
Show that your method always works.
\end{examsubparts}
\item
This problem has two parts.
\begin{examsubparts}
\item[\textnormal{(1)}]
Explain Romberg’s technique for numerical approximation of an integral.
\item[\textnormal{(2)}]
Discuss the ideas and concepts that explain why the method works.
\end{examsubparts}
\item
This problem has two parts.
\begin{examsubparts}
\item[\textnormal{(1)}]
Set up the system of linear equations to be solved when a finite difference approach is used to solve the two point boundary value problem: \(v''+b(x)v'+c(x)v=d(x)\), where \(v(0)=0\) and \(v(1)=0\).
\item[\textnormal{(2)}]
Set up the system of linear equations to be solved when a collocation approach (with basis functions \(\sin(k\pi x)\)) is used to solve the two point boundary value problem: \(v''+b(x)v'+c(x)v=d(x)\), where \(v(0)=0\) and \(v(1)=0\).
\end{examsubparts}
\item
This problem has three parts.
\begin{examsubparts}
\item[\textnormal{(1)}]
Describe the defining properties of the Legendre polynomials on the interval \([-1,1]\).
\item[\textnormal{(2)}]
Give a careful statement of the theorem fundamental to Gauss Quadrature.
\item[\textnormal{(3)}]
Show how the Gauss Quadrature method can be used to approximate the integral 
\[
\int_0^1 \exp(x^2)\,dx.
\]
\end{examsubparts}
\end{examproblems}
\end{document}
